\documentclass[UTF8]{ctexart}
\usepackage{amsmath, amssymb, geometry}
\usepackage{xcolor}
\usepackage{enumitem}
\usepackage{tcolorbox}
\usepackage{cases}
\usepackage{fancyhdr}
\geometry{a4paper, margin=1in}

% 定义红色环境
\newenvironment{redenv}{\color{red}}{}

% 定义详细解析环境
\newenvironment{detailedanalysis}{%
    \color{red}
    \begin{enumerate}[label=\textbf{第\arabic*步：}, leftmargin=*, align=left, itemsep=1.5ex]
}{%
    \end{enumerate}
}

% 定义概念框
\newenvironment{conceptbox}{%
    \begin{tcolorbox}[colback=red!5, colframe=red!50!black, title=概念解析, fonttitle=\bfseries]
}{%
    \end{tcolorbox}
}

% 定义公式推导环境
\newenvironment{formuladerivation}{%
    \begin{enumerate}[label={\textbf{公式\arabic*：}}, leftmargin=*, align=left]
}{%
    \end{enumerate}
}

% 定义题目和解析的分隔线
\newcommand{\problemrule}{\noindent\rule{\textwidth}{1.5pt}\vspace{-0.8em}}
\newcommand{\analysisrule}{\noindent\rule{\textwidth}{0.8pt}\vspace{-0.8em}}

% 宏定义：方便排版选择题选项
\newcommand{\choicesFour}[4]{
    \begin{enumerate}[label=\Alph*., itemsep=0pt, parsep=0pt, topsep=5pt, labelsep=0.5em]
        \begin{minipage}{0.22\linewidth} \item #1 \end{minipage}
        \begin{minipage}{0.22\linewidth} \item #2 \end{minipage}
        \begin{minipage}{0.22\linewidth} \item #3 \end{minipage}
        \begin{minipage}{0.22\linewidth} \item #4 \end{minipage}
    \end{enumerate}
}
\newcommand{\choicesTwo}[4]{
    \begin{enumerate}[label=\Alph*., itemsep=0pt, parsep=0pt, topsep=5pt, labelsep=0.5em]
        \begin{minipage}{0.45\linewidth} \item #1 \end{minipage}
        \begin{minipage}{0.45\linewidth} \item #2 \end{minipage}
        \begin{minipage}{0.45\linewidth} \item #3 \end{minipage}
        \begin{minipage}{0.45\linewidth} \item #4 \end{minipage}
    \end{enumerate}
}
\newcommand{\choices}[4]{
    \begin{enumerate}[label=\Alph*., itemsep=0pt, parsep=0pt, topsep=5pt, labelsep=0.5em]
        \item #1
        \item #2
        \item #3
        \item #4
    \end{enumerate}
}

\newcommand{\blank}[1]{\underline{\hspace{#1}}}

\begin{document}

\section*{第九章 多元微积分（提高）}

% ========== 第1题 ==========
\problemrule
\section*{【题目1】混合偏导数计算}
已知函数 $ z=\frac{\sin(xy)}{y} $ ，则 $ \frac{\partial^{2}z}{\partial x\partial y}= $ \blank{1cm}。
\choicesFour{$ -x\sin(xy) $}{$ x\sin(xy) $}{$ -x\cos(xy) $}{$ x\cos(xy) $}

\begin{redenv}
\analysisrule
\subsection*{逐步解析}

\begin{detailedanalysis}
\item \textbf{详细求解过程}
\begin{enumerate}[label=(\arabic*)]
\item \textbf{先求 $\partial z / \partial x$}：
$$ \frac{\partial z}{\partial x} = \frac{1}{y} \cdot \cos(xy) \cdot y = \cos(xy) $$
\item \textbf{再求 $\partial^2 z / \partial x \partial y$}：
$$ \frac{\partial}{\partial y} (\cos(xy)) = -\sin(xy) \cdot x = -x\sin(xy) $$
\end{enumerate}

\item \textbf{最终答案}
\textbf{答案}：A
\end{detailedanalysis}
\end{redenv}

% ========== 第2题 ==========
\problemrule
\section*{【题目2】可微与偏导}
二元函数 $ f(x,y) $ 在点 $ (x_{0},y_{0}) $ 处存在偏导数是 $ f(x,y) $ 在该点可微的 \blank{1cm}。
\choicesTwo
{必要而不充分条件}
{充分而不必要条件}
{必要且充分条件}
{既不必要也不充分条件}

\begin{redenv}
\analysisrule
\subsection*{逐步解析}

\begin{detailedanalysis}
\item \textbf{详细求解过程}
\begin{enumerate}[label=(\arabic*)]
\item \textbf{必要性}：可微 $\implies$ 偏导数存在。正确。
\item \textbf{充分性}：偏导数存在 $\nRightarrow$ 可微。反例 $f(x,y) = \sqrt{|xy|}$ 在 $(0,0)$。
\item \textbf{结论}：偏导数存在是可微的必要不充分条件。
\end{enumerate}

\item \textbf{最终答案}
\textbf{答案}：A
\end{detailedanalysis}
\end{redenv}

% ========== 第3题 ==========
\problemrule
\section*{【题目3】极值与驻点}
函数 $ f(x,y) $ 可微， $ f'_{x}(x_{0},y_{0})=0, f'_{y}(x_{0},y_{0})=0 $ ，则函数 $ f(x,y) $ 在 $ (x_{0},y_{0}) $ \blank{1cm}。
\choicesTwo
{可能有极值，也可能没有极值}
{必有极值，可能是极大值也可能是极小值}
{一定没有极值}
{必有极值，且为极小值}

\begin{redenv}
\analysisrule
\subsection*{逐步解析}

\begin{detailedanalysis}
\item \textbf{详细求解过程}
驻点（偏导数为0的点）不一定是极值点（例如马鞍面 $z=xy$ 在 $(0,0)$）。
需要通过二阶偏导数判定 $(AC-B^2)$。
所以是可能有，也可能没有。

\item \textbf{最终答案}
\textbf{答案}：A
\end{detailedanalysis}
\end{redenv}

% ========== 第4题 ==========
\problemrule
\section*{【题目4】二重极限}
$ \lim_{x\to0}\frac{xy^{2}}{x^{2}+y^{4}}= $ \blank{1cm}。
\choicesFour{0}{1}{$ \frac{1}{2} $}{不存在}

\begin{redenv}
\analysisrule
\subsection*{逐步解析}

\begin{detailedanalysis}
\item \textbf{详细求解过程}
\begin{enumerate}[label=(\arabic*)]
\item \textbf{路径一}：取 $x = ky^2$。
$$ \lim_{y \to 0} \frac{ky^2 \cdot y^2}{(ky^2)^2 + y^4} = \lim_{y \to 0} \frac{k y^4}{k^2 y^4 + y^4} = \frac{k}{k^2+1} $$
\item \textbf{结论}：极限值依赖于路径（斜率 $k$），所以极限不存在。
例如取 $x=y^2$ 极限为 $1/2$；取 $x=0$ 极限为 $0$。
\end{enumerate}

\item \textbf{最终答案}
\textbf{答案}：D
\end{detailedanalysis}
\end{redenv}

% ========== 第5题 ==========
\problemrule
\section*{【题目5】可微的充分条件}
$ f_{x}^{'}(x,y) $ 和 $ f_{y}^{'}(x,y) $ 在点 $ (x_{0},y_{0}) $ 连续是 $ f(x,y) $ 在点 $ (x_{0},y_{0}) $ 可微分的 \blank{1cm}。
\choicesFour{充分条件}{必要条件}{充要条件}{无关条件}

\begin{redenv}
\analysisrule
\subsection*{逐步解析}

\begin{detailedanalysis}
\item \textbf{详细求解过程}
定理：若偏导数在某点连续，则函数在该点可微。
这是充分条件。
（可微不一定偏导数连续，仅偏导数存在不一定可微）。

\item \textbf{最终答案}
\textbf{答案}：A
\end{detailedanalysis}
\end{redenv}

% ========== 第6题 ==========
\problemrule
\section*{【题目6】偏导数计算}
设 $ z = \arctan \frac{y}{x} $ ，则 $ \frac{\partial z}{\partial x} = $ \blank{2cm}。

\begin{redenv}
\analysisrule
\subsection*{逐步解析}

\begin{detailedanalysis}
\item \textbf{详细求解过程}
$$ \frac{\partial z}{\partial x} = \frac{1}{1+(\frac{y}{x})^2} \cdot \frac{\partial}{\partial x}(\frac{y}{x}) $$
$$ = \frac{x^2}{x^2+y^2} \cdot (-\frac{y}{x^2}) = -\frac{y}{x^2+y^2} $$

\item \textbf{最终答案}
\textbf{答案}：$ -\frac{y}{x^2+y^2} $
\end{detailedanalysis}
\end{redenv}

% ========== 第7题 ==========
\problemrule
\section*{【题目7】全微分形式}
设 $ z = \ln\sqrt{x^{2} + y^{2}} $ , 则 $ x\frac{\partial z}{\partial x} + y\frac{\partial z}{\partial y} = $ \blank{2cm}。

\begin{redenv}
\analysisrule
\subsection*{逐步解析}

\begin{detailedanalysis}
\item \textbf{详细求解过程}
化简：$z = \frac{1}{2} \ln(x^2+y^2)$。
\begin{enumerate}[label=(\arabic*)]
\item \textbf{求 $z_x$}：
$z_x = \frac{1}{2} \frac{2x}{x^2+y^2} = \frac{x}{x^2+y^2}$。
\item \textbf{求 $z_y$}：
$z_y = \frac{1}{2} \frac{2y}{x^2+y^2} = \frac{y}{x^2+y^2}$。
\item \textbf{代入}：
$$ x z_x + y z_y = \frac{x^2}{x^2+y^2} + \frac{y^2}{x^2+y^2} = 1 $$
\end{enumerate}

\item \textbf{最终答案}
\textbf{答案}：1
\end{detailedanalysis}
\end{redenv}

% ========== 第8题 ==========
\problemrule
\section*{【题目8】欧拉定理验证}
设 $ z = \ln(\sqrt{x} + \sqrt{y}) $ ，则 $ x \frac{\partial z}{\partial x} + y \frac{\partial z}{\partial y} = $ \blank{2cm}。

\begin{redenv}
\analysisrule
\subsection*{逐步解析}

\begin{detailedanalysis}
\item \textbf{详细求解过程}
\begin{enumerate}[label=(\arabic*)]
\item \textbf{求 $z_x$}：
$z_x = \frac{1}{\sqrt{x}+\sqrt{y}} \cdot \frac{1}{2\sqrt{x}}$。
$x z_x = \frac{\sqrt{x}}{2(\sqrt{x}+\sqrt{y})}$。
\item \textbf{求 $z_y$}：
$z_y = \frac{1}{\sqrt{x}+\sqrt{y}} \cdot \frac{1}{2\sqrt{y}}$。
$y z_y = \frac{\sqrt{y}}{2(\sqrt{x}+\sqrt{y})}$。
\item \textbf{相加}：
$$ x z_x + y z_y = \frac{\sqrt{x}+\sqrt{y}}{2(\sqrt{x}+\sqrt{y})} = \frac{1}{2} $$
\end{enumerate}

\item \textbf{最终答案}
\textbf{答案}：$ \frac{1}{2} $
\end{detailedanalysis}
\end{redenv}

% ========== 第9题 ==========
\problemrule
\section*{【题目9】求极值}
二元函数 $ f(x,y)=x^{2}(2+y^{2})+y\ln y $ 的极小值 \blank{2cm}。

\begin{redenv}
\analysisrule
\subsection*{逐步解析}

\begin{detailedanalysis}
\item \textbf{详细求解过程}
\begin{enumerate}[label=(\arabic*)]
\item \textbf{求一阶偏导}：
$f_x = 2x(2+y^2)$。
$f_y = 2x^2 y + \ln y + 1$。
\item \textbf{解驻点}：
$f_x = 0 \implies x=0$（因为 $2+y^2 \ge 2 \neq 0$）。
代入 $f_y=0$：$\ln y + 1 = 0 \implies \ln y = -1 \implies y = 1/e$。
驻点 $(0, 1/e)$。
\item \textbf{二阶偏导判别}：
$f_{xx} = 2(2+y^2)$。
$f_{xy} = 4xy$。
$f_{yy} = 2x^2 + 1/y$。
在 $(0, 1/e)$ 处：
$A = f_{xx} = 2(2 + e^{-2}) > 0$。
$B = f_{xy} = 0$。
$C = f_{yy} = e$。
$AC - B^2 = 2e(2+e^{-2}) > 0$。
$A > 0$，所以是极小值点。
\item \textbf{求极小值}：
$f(0, 1/e) = 0 + \frac{1}{e} \ln(\frac{1}{e}) = -\frac{1}{e}$。
\end{enumerate}

\item \textbf{最终答案}
\textbf{答案}：$ -\frac{1}{e} $
\end{detailedanalysis}
\end{redenv}

% ========== 第10题 ==========
\problemrule
\section*{【题目10】定义域求解}
$ z = \ln(x^{2} + y^{2} - 2) + \sqrt{4 - x^{2} - y^{2}} $ 函数的定义域为 \blank{2cm}。

\begin{redenv}
\analysisrule
\subsection*{逐步解析}

\begin{detailedanalysis}
\item \textbf{详细求解过程}
条件：
1. $x^2 + y^2 - 2 > 0 \implies x^2 + y^2 > 2$。
2. $4 - x^2 - y^2 \ge 0 \implies x^2 + y^2 \le 4$。
综合：$2 < x^2 + y^2 \le 4$。
几何意义：圆环区域（不含内边界）。

\item \textbf{最终答案}
\textbf{答案}：$ 2 < x^2 + y^2 \le 4 $
\end{detailedanalysis}
\end{redenv}

% ========== 第11题 ==========
\problemrule
\section*{【题目11】复合函数偏导}
已知函数 $ z = f(u, v) $ 可导， $ u = x \arcsin y, v = \frac{y}{x} $ ，求 $ \frac{\partial z}{\partial x}, \frac{\partial z}{\partial y} $。

\begin{redenv}
\analysisrule
\subsection*{逐步解析}

\begin{detailedanalysis}
\item \textbf{详细求解过程}
利用链式法则。
\begin{enumerate}[label=(\arabic*)]
\item \textbf{求 $\partial z / \partial x$}：
$$ \frac{\partial z}{\partial x} = f_u \frac{\partial u}{\partial x} + f_v \frac{\partial v}{\partial x} $$
$$ = f_u (\arcsin y) + f_v (-\frac{y}{x^2}) = \arcsin y \cdot f_u - \frac{y}{x^2} f_v $$
\item \textbf{求 $\partial z / \partial y$}：
$$ \frac{\partial z}{\partial y} = f_u \frac{\partial u}{\partial y} + f_v \frac{\partial v}{\partial y} $$
$$ = f_u (\frac{x}{\sqrt{1-y^2}}) + f_v (\frac{1}{x}) = \frac{x}{\sqrt{1-y^2}} f_u + \frac{1}{x} f_v $$
\end{enumerate}

\item \textbf{最终答案}
\textbf{答案}：
$ \frac{\partial z}{\partial x} = \arcsin y f_{u} - \frac{y}{x^{2}} f_{v} $;
$ \frac{\partial z}{\partial y} = \frac{x}{\sqrt{1-y^{2}}} f_{u} + \frac{1}{x} f_{v} $
\end{detailedanalysis}
\end{redenv}

% ========== 第12题 ==========
\problemrule
\section*{【题目12】拉格朗日乘数法应用}
求内接于半径为 $R$ 的球体且有最大体积的长方体的边长。

\begin{redenv}
\analysisrule
\subsection*{逐步解析}

\begin{detailedanalysis}
\item \textbf{详细求解过程}
设长方体边长为 $2x, 2y, 2z$（为了对称方便）。
顶点均在球面上：$x^2 + y^2 + z^2 = R^2$。
体积 $V = 8xyz$。需要最大化 $xyz$。
利用均值不等式：
$$ \sqrt[3]{x^2 y^2 z^2} \le \frac{x^2+y^2+z^2}{3} = \frac{R^2}{3} $$
当且仅当 $x^2 = y^2 = z^2$ 时取等号。
即 $x=y=z = \frac{R}{\sqrt{3}}$。
此时边长为 $\frac{2R}{\sqrt{3}}$。
长方体为正方体。
边长 $a = b = c = \frac{2\sqrt{3}}{3}R$。

\item \textbf{最终答案}
\textbf{答案}：$ \frac{2\sqrt{3}}{3}R $
\end{detailedanalysis}
\end{redenv}

% ========== 第13题 ==========
\problemrule
\section*{【题目13】隐函数求导}
设函数 $ z = z(x, y) $ 由方程 $ z^{3} - 3xyz = a^{3} $ 确定，求 $ \frac{\partial z}{\partial x}, \frac{\partial z}{\partial y} $ 和 $dz$。

\begin{redenv}
\analysisrule
\subsection*{逐步解析}

\begin{detailedanalysis}
\item \textbf{详细求解过程}
设 $F(x,y,z) = z^3 - 3xyz - a^3 = 0$。
\begin{enumerate}[label=(\arabic*)]
\item \textbf{求偏导}：
$F_x = -3yz$。
$F_y = -3xz$。
$F_z = 3z^2 - 3xy$。
\item \textbf{公式法}：
$$ \frac{\partial z}{\partial x} = -\frac{F_x}{F_z} = -\frac{-3yz}{3z^2 - 3xy} = \frac{yz}{z^2-xy} $$
$$ \frac{\partial z}{\partial y} = -\frac{F_y}{F_z} = -\frac{-3xz}{3z^2 - 3xy} = \frac{xz}{z^2-xy} $$
\item \textbf{求全微分 dz}：
$$ dz = \frac{\partial z}{\partial x} dx + \frac{\partial z}{\partial y} dy = \frac{z}{z^2-xy}(y dx + x dy) $$
\end{enumerate}

\item \textbf{最终答案}
\textbf{答案}：
$ \frac{\partial z}{\partial x} = \frac{yz}{z^2-xy}, \frac{\partial z}{\partial y} = \frac{xz}{z^2-xy} $;
$ dz = \frac{z}{z^2-xy}(y dx + x dy) $
\end{detailedanalysis}
\end{redenv}

% ========== 第14题 ==========
\problemrule
\section*{【题目14】高阶偏导数}
设 $ z = f(2x - y) + g(x, xy) $ ，其中函数 $ f(w) $ 具有二阶导数， $ g(u, v) $ 具有二阶连续偏导数，求 $ \frac{\partial z}{\partial x} $ 与 $ \frac{\partial^{2}z}{\partial x \partial y} $。

\begin{redenv}
\analysisrule
\subsection*{逐步解析}

\begin{detailedanalysis}
\item \textbf{详细求解过程}
\begin{enumerate}[label=(\arabic*)]
\item \textbf{求 $\frac{\partial z}{\partial x}$}：
对 $f(2x-y)$：$f' \cdot 2 = 2f'$。
对 $g(x, xy)$：$g_1 \cdot 1 + g_2 \cdot y = g_1 + y g_2$。
所以 $\frac{\partial z}{\partial x} = 2f'(2x-y) + g_1(x, xy) + y g_2(x, xy)$。
\item \textbf{求 $\frac{\partial^2 z}{\partial x \partial y}$}：
对 $2f'(2x-y)$ 求 $y$ 偏导：$2f'' \cdot (-1) = -2f''$。
对 $g_1(x, xy)$ 求 $y$ 偏导：$g_{12} \cdot x$。
对 $y g_2(x, xy)$ 求 $y$ 偏导：$1 \cdot g_2 + y(g_{22} \cdot x) = g_2 + xy g_{22}$。
综合：
$$ \frac{\partial^2 z}{\partial x \partial y} = -2f'' + x g_{12} + g_2 + xy g_{22} $$
\end{enumerate}

\item \textbf{最终答案}
\textbf{答案}：
$ \frac{\partial z}{\partial x} = 2f' + g_1 + y g_2 $;
$ \frac{\partial^{2}z}{\partial x \partial y} = -2f'' + x g_{12} + g_2 + xy g_{22} $
\end{detailedanalysis}
\end{redenv}

% ========== 第15题 ==========
\problemrule
\section*{【题目15】条件极值几何应用}
求内接于椭圆 $ \frac{x^{2}}{3}+\frac{y^{2}}{4}=1 $ 的最大长方形的面积。

\begin{redenv}
\analysisrule
\subsection*{逐步解析}

\begin{detailedanalysis}
\item \textbf{详细求解过程}
设长方形在第一象限顶点为 $(x,y)$。面积 $S = 4xy$。
条件：$\frac{x^2}{3} + \frac{y^2}{4} = 1$。
令 $x = \sqrt{3}\cos t, y = 2\sin t$ ($t \in (0, \pi/2)$)。
$S = 4\sqrt{3}\cos t \cdot 2\sin t = 4\sqrt{3} \sin 2t$。
当 $\sin 2t = 1$ 即 $t = \pi/4$ 时面积最大。
最大面积 $S_{max} = 4\sqrt{3}$。
（也可以用拉格朗日乘数法或均值不等式 $\frac{x^2}{3} \cdot \frac{y^2}{4} \le (\frac{1}{2})^2$）。

\item \textbf{最终答案}
\textbf{答案}：$ 4\sqrt{3} $
\end{detailedanalysis}
\end{redenv}

\end{document}
